% ============================================================
% Model: NaCl 几何结构因子与 X 射线衍射
% ============================================================

\section{NaCl 几何结构因子与 X 射线衍射}
\label{sec:xray-nacl}

\begin{tipbox}[关键词标签]
几何结构因子 · X射线衍射 · NaCl结构 · 消光规则 · 布拉格定律 · 面心立方
\end{tipbox}

% --------------------------------------------------------
% 1. 模型定义框
% --------------------------------------------------------
\begin{modelbox}[模型：NaCl 结构的 X 射线衍射]
\textbf{物理情境}：X 射线衍射是测定晶体结构的最重要实验手段。衍射强度由几何结构因子 $S(\mathbf{G})$ 决定，而 $S(\mathbf{G})$ 反映了晶胞内原子的种类和位置。NaCl 结构是最典型的离子晶体结构之一。

\textbf{NaCl 结构描述}：
\begin{itemize}[nosep]
    \item 布拉维格子：面心立方（FCC）
    \item 基元：Na$^+$ 在 $(0,0,0)$，Cl$^-$ 在 $(1/2,1/2,1/2)$（以 FCC 格点坐标）
    \item 每个晶胞含 4 个 Na$^+$ 和 4 个 Cl$^-$
\end{itemize}

\textbf{关键公式}：
\begin{itemize}[nosep]
    \item 布拉格定律：$2d\sin\theta=\lambda$
    \item 劳厄条件：$\mathbf{k}'-\mathbf{k}=\mathbf{G}$
    \item 几何结构因子：$S(\mathbf{G})=\sum_j f_j\,e^{i\mathbf{G}\cdot\mathbf{d}_j}$
\end{itemize}
\end{modelbox}

% --------------------------------------------------------
% 2. 关键公式框
% --------------------------------------------------------
\begin{formulabox}[核心公式]
\begin{enumerate}[label=\textbf{(\arabic*)}]
    \item \textbf{FCC 倒格矢}：
    \[
    \mathbf{G}=h\mathbf{b}_1+k\mathbf{b}_2+l\mathbf{b}_3=\frac{2\pi}{a}(h,k,l)
    \]
    \texteq{FCC 的衍射选择定则：$h,k,l$ 全奇或全偶才出现衍射}

    \item \textbf{NaCl 几何结构因子}：
    \[
    S_{hkl}=f_++f_-\,e^{i\pi(h+k+l)}
    =\begin{cases}
    f_++f_-, & h+k+l\text{ 为偶数},\\
    f_+-f_-, & h+k+l\text{ 为奇数}.
    \end{cases}
    \]
    \texteq{$f_+$ 为正离子散射因子，$f_-$ 为负离子散射因子}

    \item \textbf{特殊情形 ($f_+=f_-=f$)}：
    \[
    S_{hkl}=\begin{cases}
    2f, & h+k+l\text{ 为偶数},\\
    0, & h+k+l\text{ 为奇数}.
    \end{cases}
    \]
    \texteq{奇数混合指数衍射峰完全消失（如 $(111), (311)$ 等）}

    \item \textbf{衍射强度}：
    \[
    I_{hkl}\propto|S_{hkl}|^2\cdot|F_{\text{FCC}}|^2
    \]
    \texteq{$F_{\text{FCC}}$ 为 FCC 格子的几何因子，决定全奇/全偶选择定则}

    \item \textbf{晶面间距}：
    \[
    d_{hkl}=\frac{a}{\sqrt{h^2+k^2+l^2}}
    \]
\end{enumerate}
\end{formulabox}

% --------------------------------------------------------
% 3. 训练点框
% --------------------------------------------------------
\begin{drillbox}[训练点与考法变体]
\trainingpoint{1} \textbf{结构因子直接计算}：给定晶体结构（NaCl、CsCl、金刚石、ZnS 等），写出基元坐标，计算 $S_{hkl}$，推导消光规则。

\trainingpoint{2} \textbf{消光规则应用}：给定衍射图谱，根据出现/缺失的衍射峰指数，反推晶体结构类型（如 $(111)$ 峰缺失说明 $f_+=f_-$）。

\trainingpoint{3} \textbf{结构因子与原子散射因子}：$f_j$ 随散射角增大而减小（原子形状因子），导致高角度衍射峰强度降低。

\trainingpoint{4} \textbf{温度因子}：原子热振动使衍射峰强度降低，引入 Debye-Waller 因子 $e^{-2W}$，$W\propto T$。

\trainingpoint{5} \textbf{与电子衍射/中子衍射对比}：X 射线被电子云散射；电子被库仑势散射（强得多）；中子被原子核散射，对轻元素敏感。
\end{drillbox}

% --------------------------------------------------------
% 4. 解题技巧框
% --------------------------------------------------------
\begin{tipbox}[快速识别与解题技巧]
\begin{itemize}
    \item \examnote{识别标志}：题干出现``结构因子''''消光规则''''X射线衍射''''布拉格定律''，立即定位本模型。
    \item \examnote{常见结构速查}：
    \begin{center}
    \begin{tabular}{lll}
    \toprule
    \textbf{结构} & \textbf{基元} & \textbf{特殊消光} \\
    \midrule
    SC & 单原子 & 无 \\
    BCC & 体心 & $h+k+l$=奇数消光 \\
    FCC & 面心 & $h,k,l$ 奇偶混合消光 \\
    NaCl & FCC+双原子 & 同上，且 $f_+=f_-$ 时奇数和消光 \\
    CsCl & SC+双原子 & $h+k+l$=奇数时 $S=f_+-f_-$ \\
    金刚石 & FCC+双原子 & $h+k+l=4n+2$ 时消光 \\
    \bottomrule
    \end{tabular}
    \end{center}
    \item \examnote{FCC 选择定则记忆}：FCC 的消光来源于面心原子引起的相消干涉。若 $h,k,l$ 奇偶混合，则四个 FCC 原子贡献相互抵消。
    \item \examnote{NaCl 的特殊消光}：当 $f_+=f_-$ 时，NaCl 的奇数混合指数峰消失，此时衍射图样看起来像简单立方（但晶格常数为实际的一半）。这是因为正负离子散射完全相消。
    \item \examnote{结构因子是复数}：$S=|S|e^{i\phi}$，衍射强度只与 $|S|^2$ 有关，但结构因子本身的相位对结构解析（如直接法、Patterson 法）至关重要。
\end{itemize}
\end{tipbox}

% --------------------------------------------------------
% 5. 易错点框
% --------------------------------------------------------
\begin{errorbox}{常见失误与陷阱}
\begin{itemize}
    \item \textbf{基矢坐标系}：FCC 的基元坐标常用分数坐标（以立方胞边长 $a$ 为单位），如 $(0,0,0)$、$(1/2,1/2,0)$ 等。不要误用实际长度坐标。
    \item \textbf{全奇与全偶}：FCC 的衍射条件是 $h,k,l$ \textbf{全奇}或\textbf{全偶}，不是 $h+k+l$=偶数。例如 $(200)$ 允许，但 $(100)$ 不允许（虽然 $1+0+0$=奇数，但关键是混合）。
    \item \textbf{NaCl 与金刚石的区别}：两者都是 FCC+双原子，但基元位置不同。NaCl：$(0,0,0)$ 和 $(1/2,1/2,1/2)$；金刚石：$(0,0,0)$ 和 $(1/4,1/4,1/4)$。金刚石的消光更严格（$h+k+l$ 必须是 $4n$ 或全奇）。
    \item \textbf{消光不等于零强度}：严格消光时 $S=0$，强度为零。但若 $f_+\neq f_-$（如 NaCl），``额外''的消光条件（$h+k+l$=奇数时 $S=f_+-f_-$）不意味着完全消光，只是强度较弱。
    \item \textbf{倒格矢与晶面指数}：$\mathbf{G}_{hkl}$ 垂直于 $(hkl)$ 晶面，$|\mathbf{G}_{hkl}|=2\pi/d_{hkl}$。布拉格定律等价于劳厄条件。
\end{itemize}
\end{errorbox}

% --------------------------------------------------------
% 6. 物理图像与关联模型
% --------------------------------------------------------
\begin{insightbox}[物理图像与模型关联]
\textbf{物理图像}：

X 射线衍射就像用光照射一个周期性排列的栅栏：
\begin{itemize}[nosep]
    \item \textbf{布拉格定律}：只有当入射角满足``反射''条件时，来自不同``栅栏条''（晶面）的散射波才会相长干涉。
    \item \textbf{结构因子}：每个``栅栏条''内部有多个小孔（晶胞内的原子），小孔的相对位置和大小决定了散射波的相对相位和振幅。如果小孔排列使得散射波相互抵消，就出现``消光''。
    \item \textbf{NaCl 的双原子基元}：每个栅栏条上有两种大小不同的小孔（Na$^+$ 和 Cl$^-$），它们的散射波在某些角度相长，在另一些角度相消。
\end{itemize}

\textbf{与相似模型的对比}：
\begin{center}
\begin{tabular}{llll}
\toprule
\textbf{对比维度} & \textbf{X射线衍射} & \textbf{电子衍射} & \textbf{中子衍射} \\
\midrule
散射源 & 电子云 & 库仑势（整体） & 原子核 \\
散射强度 & 弱（$\sim10^{-4}$） & 强 & 中等 \\
对轻元素 & 不敏感 & 敏感 & 非常敏感 \\
磁结构 & 不可见 & 不可见 & 可见（磁散射） \\
\bottomrule
\end{tabular}
\end{center}

\textbf{推广方向}：
\begin{itemize}[nosep]
    \item 反常散射：利用吸收边附近的波长变化确定原子位置（如蛋白质晶体学中的 MAD 法）
    \item 表面衍射：低能电子衍射（LEED）测定表面原子结构
    \item 时间分辨 X 射线衍射：超快激光泵浦探测晶格动力学
\end{itemize}
\end{insightbox}

% --------------------------------------------------------
% 7. 推导附录
% --------------------------------------------------------
\begin{derivationbox}[推导细节（自我检验用）]
\textbf{Step 1：FCC 基元坐标}

以原点在顶点，FCC 的四个原子位置：
\[
\mathbf{d}_1=(0,0,0),\;\mathbf{d}_2=(\tfrac{1}{2},\tfrac{1}{2},0),\;
\mathbf{d}_3=(\tfrac{1}{2},0,\tfrac{1}{2}),\;\mathbf{d}_4=(0,\tfrac{1}{2},\tfrac{1}{2}).
\]

\textbf{Step 2：FCC 结构因子}

\begin{align*}
F_{\text{FCC}}&=f\bigl[1+e^{i\pi(h+k)}+e^{i\pi(h+l)}+e^{i\pi(k+l)}\bigr]\\
&=f\bigl[1+(-1)^{h+k}+(-1)^{h+l}+(-1)^{k+l}\bigr].
\end{align*}

若 $h,k,l$ 全偶或全奇：每一项都是 $+1$，$F_{\text{FCC}}=4f$。

若 $h,k,l$ 奇偶混合（如两偶一奇）：两项为 $+1$，两项为 $-1$，$F_{\text{FCC}}=0$。

\textbf{Step 3：NaCl 结构因子}

NaCl 可视为 FCC 格子，格点上交替放 Na$^+$ 和 Cl$^-$。等价描述：FCC 格子 + 基元（Na$^+$ 在 0，Cl$^-$ 在 $(1/2,1/2,1/2)$）。

\[
S_{hkl}=f_+\,F_{\text{FCC}}+f_-\,F_{\text{FCC}}\,e^{i\pi(h+k+l)}
=F_{\text{FCC}}\bigl[f_++f_-(-1)^{h+k+l}\bigr].
\]

若 $h,k,l$ 奇偶混合：$F_{\text{FCC}}=0$，$S=0$（FCC 消光）。

若 $h,k,l$ 全偶：$S=4(f_++f_-)$。

若 $h,k,l$ 全奇：$S=4(f_+-f_-)$。

\textbf{Step 4：$f_+=f_-$ 的特殊情形}

当正负离子散射因子相等（如 X 射线对核电荷数相近的元素），全奇衍射峰 $S=0$，只保留全偶峰。此时衍射图样类似简单立方，晶格常数为 $a/2$。
\end{derivationbox}

\vspace{1em}
